\begin{abstract}
\noindent
The Zenodo record 10.5281/zenodo.22678442, \emph{Verification and Reconstruction of a Claimed Complex Structure on $S^6$: An AI-Produced Validate/Repair/Disprove Record}, collects the owner's AI-produced audit of a manuscript circulated by Levent Alpöge. The manuscript claims that a family of complex two-dimensional tori over the projective line, completed at three special fibres, is a complex structure on the six-sphere. The record holds 33 files from four versions (1--9 September 2026): a 154-page verification monograph, the reconstructed and annotated manuscript, an eleven-page correction note on a step in a published theorem of Campana, Demailly and Peternell, topology supplements, and a 174~MB archive of the full project.

This reader maps the record, states what it claims, with its own labels, and says how far each claim was checked here. The central claims are:
\begin{itemize}
\item the construction gives a compact complex threefold diffeomorphic to $S^6$, of algebraic dimension one;
\item it is therefore a counterexample to the published theorem;
\item a specific step of that theorem's proof is not a formal consequence of its displayed sequences. The source manuscript had already made this point for its own fibre; the record adds an independent test family.
\end{itemize}
The first two are claims of the record and of the source manuscript. An independent exposition by Philip Engel (13 September 2026) presents the construction with the stated aim of a self-contained proof. Neither the construction nor the counterexample claim was checked here.

What this reader does check (Section~\ref{sec:correction}):
\begin{itemize}
\item the correction note, in full. Its test family shows that the step in question needs an extra hypothesis. Added here:
\begin{itemize}
\item the step is valid on normal fibres, as the source manuscript already notes, and the torsion that breaks it is present at every non-normal point;
\item at a reduced fibre with normal crossings, the group that the step asserts to vanish is the space of sections of the line bundle pulled back to the normalization. So the step fails exactly when that bundle has a nonzero section;
\item in the test family this group is one-dimensional for every gluing parameter;
\end{itemize}
\item the local algebra behind one of the record's corrections to the source manuscript, in general form. The correction is right for relative one-forms, but relative two-forms do have torsion at the cusp fibre;
\item a number of finite statements, and the exact set of values of the cubic in the key advance S6-5.
\end{itemize}
This version includes the findings of one independent referee pass (Appendix~\ref{app:verification}).
\end{abstract}

\section*{About this reader}\phantomsection\label{sec:about}
\addcontentsline{toc}{section}{About this reader}

\paragraph{The record, and who made it.}
\begin{itemize}
\item The candidate construction and its original claims are attributed by the record to their reported source. Different files name the source differently:
\begin{itemize}
\item the monograph's title page says ``Claude/Fabel'';
\item the key-advances reader of 6 September says the construction was circulated by Levent Alpöge with Fable;
\item the later attribution file of the workbench (\texttt{ATTRIBUTION.md} of the Yang--Mills repository, 21 September), which its author designates as the clarification, says the manuscript was circulated by Levent Alpöge and produced with Claude.
\end{itemize}
\item The verification monograph describes itself as Codex's mathematical verification, with two independent ChatGPT Pro review records. The record calls itself an AI-produced verification aid, not specialist certification, and makes no priority claim (record \texttt{README.md}).
\item The record's creator, in its metadata, is KokunoYumeto. Its provenance file records that the user of the underlying Codex session directed that the programme be called validate/repair/disprove, and that it put proofs, calculations and counterexamples before verdicts.
\item The correction note's author line reads ``human author to be supplied''.
\end{itemize}

\paragraph{What was read.}
\begin{itemize}
\item The record's files, downloaded on 26 September 2026 with the owner's permission. Their MD5 checksums equal the record's; the 174~MB archive \texttt{28\_} was not downloaded.
\item The correction note, in full.
\item From the monograph:
\begin{itemize}
\item its orientation (Part~I), including its 21-item status list;
\item its Theorem~13.6;
\item its closing dispositions (§§13.11 and 14).
\end{itemize}
Its Theorems~13.5, 13.25 and 13.28 are cited by the numbers the record gives them.
\item From the manuscript: its opening summary, its Remarks~1.3 and~9.21, and its §§10.1--10.3 (the quotation of the published step, Lemma~10.3, Remark~10.4, and the statement and first steps of Theorem~10.5).
\item The claim ledger in the record's evidence archive \texttt{04\_}.
\item The five-page key-advances reader, the reader's guide to the 5 September supplement, the guide to the frozen project, and the provenance file.
\end{itemize}
The rest was mapped from its tables of contents and summaries.

\paragraph{How it was checked, and the labels.}
The labels are those of the companion readers, with one addition:
\begin{itemize}
\item \textsf{Proved here}: the proof is given in full in this reader;
\item \textsf{Audited}: the record's proof was read in full, every step was checked, and independent checks were run;
\item \textsf{Reproduced}: the finite content was recomputed;
\item \textsf{Located}: the statement and its proof were found and are described; nothing was checked;
\item \textsf{Stated}: the statement and the location of its proof are reported, but the proof was not seen.
\end{itemize}
The record's own labels are \textsf{VERIFIED}, \textsf{OPEN} and \textsf{REFUTED}, assigned by its AI verification. They are quoted as the record's and are not endorsed here. This version includes the findings of one independent referee pass, each verified before it was applied (Appendix~\ref{app:verification}).
